\begin{ProgramList}

第\ref{Chap:chap3}章 Swift拥塞控制算法Python实现（swift-tcp文件夹）

tcp\_swift.py：Swift启发式决策模块核心实现（约380行），包含TcpSwift类，实现多信号融合拥塞检测、差异化拥塞响应、启发式自适应参数调优（$\alpha$动态调整）、融合控制策略（BDP估计与窗口调整）、每流独立状态管理等全部核心算法逻辑

tcp\_base.py：TCP状态观测基类（约140行），实现TcpEventBased（11维事件驱动）和TcpTimeBased（16维时间驱动）两种状态解析模式，负责将ns-3传递的原始状态向量解析为语义化的决策参数

test\_swift.py：Swift算法测试与运行脚本（约120行），基于ns3-gym框架的跨进程运行入口，创建状态—决策传输环境实例，为每个TCP套接字（socketUuid）创建独立的Swift决策模块实例，实现多流并发控制
\\

第\ref{Chap:chap3}章 ns-3网络仿真C++模块（swift-tcp文件夹）

tcp-swift.cc / tcp-swift.h：TcpSwift拥塞控制类实现，负责在ns-3 TCP协议栈中注册Swift算法并创建环境实例，采集协议栈状态并应用窗口决策

tcp-swift-env.cc / tcp-swift-env.h：Swift环境类（约230行），继承自ns3-gym的TcpGymEnv，实现11维状态空间的数据采集与传递、ECN感知的反馈值计算（ECN惩罚因子）、延迟拥塞窗口应用机制（Deferred Cwnd）
\\

第\ref{Chap:chap4}章 网络仿真实验配置（swift-tcp文件夹）

sim.cc：ns-3网络仿真主程序（约400行），构建哑铃型（Dumbbell）网络拓扑、配置TCP协议参数与RED/ECN队列、设置FlowMonitor流量监控、配置跨进程状态传输接口，支持通过命令行参数灵活配置36种实验场景与UDP突发干扰

Makefile：自动化编译与实验运行脚本，支持多场景批量实验、对比测试和日志收集

\end{ProgramList}
